\begin{trapbox}

	\begin{description}[style=nextline, leftmargin=2.8em, labelsep=0.5em, itemsep=0.35em]
		\item[\textbf{\textcolor{CTrap}{易错点一（共线前提）}}]
			在使用向量模长或数量积判定三角形形状时，未验证 $A,B,C$ 三点是否共线，直接当作三角形处理。
		\item[\textbf{\textcolor{CTrap}{正确理解（条件检查）}}]
			三角形定义要求三点不共线，因此必须先确认三点不共线，或题目已默认为三角形。
		\item[\textbf{\textcolor{CTrap}{错因分析（跳过检查）}}]
			学生常忽略这一前提，将共线情况当成三角形，从而误用结论。
	\end{description}

	\medskip
	\begin{description}[style=nextline, leftmargin=2.8em, labelsep=0.5em, itemsep=0.35em]
		\item[\textbf{\textcolor{CTrap}{易错点二（共起点）}}]
			误认为 $\overrightarrow{AB}\cdot\overrightarrow{BC}=0$ 即可判定 $\angle B=90^\circ$，但 $\overrightarrow{AB}$ 与 $\overrightarrow{BC}$ 起点不同，其点乘结果并不代表 $\angle B$。
		\item[\textbf{\textcolor{CTrap}{正确理解（向量夹角定义）}}]
			数量积为零判定直角时，两个向量必须起点相同，才能表示该点处的内角。
		\item[\textbf{\textcolor{CTrap}{错因分析（起点错误）}}]
			混淆向量夹角的定义，使用了终点到起点的向量组合。
	\end{description}

	\medskip
	\begin{description}[style=nextline, leftmargin=2.8em, labelsep=0.5em, itemsep=0.35em]
		\item[\textbf{\textcolor{CTrap}{易错点三（充分条件）}}]
			把 $|\overrightarrow{AB}|=|\overrightarrow{AC}|$ 当作三角形为等腰的充要条件，认为等腰三角形一定满足 $AB=AC$。
		\item[\textbf{\textcolor{CTrap}{正确理解（等腰定义）}}]
			三角形只要有两条边相等即为等腰，$AB=AC$ 只是其中一种情况，不能排除 $AB=BC$ 或 $AC=BC$ 的情况。
		\item[\textbf{\textcolor{CTrap}{错因分析（充要混淆）}}]
			受思维定势影响，默认等腰三角形有两条边相等，但未意识到相等的边不一定正好是 $AB$ 和 $AC$。
	\end{description}

\end{trapbox}
